Figure~\ref{fig:fig-two-spellings} writes one intent in both systems: the app key
may call \code{guess} on the game, and may send XLM only to the game, at most
\fact{preview-game-cap-xlm}~XLM per \fact{preview-game-period-ledgers} ledgers
(\fact{preview-game-period-label}), through ledger $N-1$. It is the game grant the
testnet preview of the \fact{wallet-name} wallet installed through Perch
(Perch \srcpath{release_stack.rs:4019-4030}). oz-policy-builder's side is an Intent
(\srcpath{packages/policy-synth/src/intent.ts:18-59}), the predicate it lowers to
for the XLM rule (\srcpath{intent.ts:117-140}) and the \code{install\_policy}
arguments (\srcpath{src/tools.ts:588-640}), read from source; we did not run these
installs. The marks in the figure match the notes.

\perchfigure{fig-two-spellings}

\begin{enumerate}[label=\protect\callout{\arabic*},leftmargin=*,itemsep=1pt]
\item Both select one rule per context, and each transfer rule sees only the
transfer (\srcpath{policy-interpreter/README.md:647-649}), so in both a transfer to
the game is allowed without a \code{guess}.
\item Perch's window cap is part of the rule; a document cannot bound one call's
amount. oz-policy-builder requires a per-call bound (\code{lte}) unless the owners
set \code{allowUnboundedAmount}, and attaches \code{spending\_limit} for the
window, because the interpreter keeps no totals (\srcpath{dsl.rs:65-72}).
\item Perch installs both rules with one \code{apply\_grant} call;
oz-policy-builder builds one \code{add\_context\_rule} per \code{install\_policy}
call (\srcpath{tools.ts:588-590}).
\item Perch's rule holds a passkey as an external signer. oz-policy-builder writes
every rule signer as delegated (\srcpath{install.ts:148}) and the interpreter refuses
external signers (error \fact{opb-error-external-signer}), so its app signs as a
Stellar or contract address; its install is authorized by \code{sourceAccount}, a
signer of rule~0 (\srcpath{tools.ts:590}). Any account may submit a Perch install
that carries the owner's authorization.
\item In Perch, \fact{perch-expiry-semantics}. In OpenZeppelin a rule is
\fact{oz-valid-until-semantics} (\srcpath{storage.rs:321-322} in
\code{stellar-contracts}); \code{validUntilLedger}~$N-1$ therefore names the same
last ledger.
\end{enumerate}

\noindent One more difference has no mark. oz-policy-builder can consult live
state, for example a \code{read} asking a registry whether the recipient is listed
(\srcpath{docs/stellar.md:369-375}); a Perch rule reads nothing outside the account.
